% !TeX program = xelatex
% ============================================================
% 论文主文件 —— paper-pipeline 阶段 7 产出
% 默认 ctexart + xelatex:中英混排都没问题(占位标记"待补实验"含中文,
% pdflatex 编不了),纯英文内容渲染与 article 无异。
% 投稿换 venue 模板时,英文 venue 可换回 \documentclass{article} + pdflatex,
% 但需把 \todobox 的中文提示语一并改英文。
% ============================================================
\documentclass[11pt,a4paper]{ctexart}

\usepackage[margin=1in]{geometry}
\usepackage{graphicx}
\usepackage{booktabs}
\usepackage{amsmath,amssymb}
\usepackage{multirow}
\usepackage[numbers,sort&compress]{natbib}
\usepackage[colorlinks=true,linkcolor=blue,citecolor=blue,urlcolor=blue]{hyperref}
\usepackage{xcolor}

% ---- 占位标注块:auto 模式下未经实验验证的数字一律用它显式标出 ----
\newcommand{\todobox}[1]{%
  \textcolor{orange}{$\bigstar$}\textcolor{orange}{\textsf{[#1]}}}

\graphicspath{{figures/}{../figures/}}

\title{<论文标题:阶段 1 的 topic 一句话版提炼,按 anti-defensive-writing
  发布会原则突出最值得发表的点>}

\author{<作者与单位;投稿版按目标期刊要求拆到 title page>}
\date{}

\begin{document}

\maketitle

% ============================================================
% 分节契约(与 ../outline.md 对齐,起草前先填全契约再动笔)
% purpose:      一段讲清问题、方法、核心结果与意义
% allowed claims: 仅 evidence_table 中状态非"待补"的 claim + 数字结论
% forbidden claims: 未登记 claim;占位数字不得以确定语气出现
% inputs:       evidence_table.md · argument_map.md
% 验证方式:     nature-ref-verifier + nature-polishing
% ============================================================
\begin{abstract}
<问题一句话。><方法核心思想一两句。><主结果:\todobox{待补实验} 数字到位后回填。>
<意义一句:该结果改变了什么。>
\end{abstract}

% ============================================================
% 分节契约
% purpose:      建立"为什么这个问题重要、为什么现有方案不够、本文做了什么"
% allowed claims: 版图性陈述须有文献行支撑(01-literature 精读卡)
% forbidden claims: 贬低竞品的无证据表述;防御性让步(anti-defensive-writing)
% inputs:       01-topic.md · 00-canonical/evidence_table.md
% 验证方式:     nature-citation 配引用 + ref-verifier 核对
% ============================================================
\section{Introduction}\label{sec:intro}

% 叙事顺序按 anti-defensive-writing 发布会原则:
% 第 1 段问题与利害 → 第 2 段现状与缺口(引文)→ 第 3 段本文做法
% → 第 4 段贡献清单 → 第 5 段结果预告
<第 1 段:领域问题与利害。>

<第 2 段:现有方法路线与缺口,配 \cite{...}。>

<第 3 段:本文方法核心思想。>

Our contributions are as follows:
\begin{itemize}
  \item <贡献 1:与 argument_map 的主 claim 一一对应>
  \item <贡献 2>
  \item <贡献 3>
\end{itemize}

% ============================================================
% 分节契约
% purpose:      给出"已解决 / 有争议 / 没人做"三分版图,定位本文
% allowed claims: 每条对已有工作的概括须能指回精读卡与 evidence_table
% forbidden claims: 只引不评的流水账;为衬托自己曲解竞品
% inputs:       01-literature/cards/*.md · evidence_table.md
% 验证方式:     nature-citation + nature-ref-verifier 逐条核对
% ============================================================
\section{Related Work}\label{sec:related}

% 按主题分小节组织(不是按时间罗列);每小节末尾一句"与本文的差异"
\subsection{<主题 A(如:消息传递框架)>}
<概述与代表工作 \cite{...}。与本文的差异:……>

\subsection{<主题 B(如:几何/三维方法)>}
<概述与代表工作 \cite{...}。与本文的差异:……>

\subsection{<主题 C(如:数据基准与评估协议)>}
<概述与代表工作 \cite{...}。与本文的差异:……>

% ============================================================
% 分节契约
% purpose:      方法形式化描述,审稿人能据此复现
% allowed claims: 每个设计选择配一句"为什么不用 X"(design.md 差异表)
% forbidden claims: 无消融支撑的优越性断言(留给 Results)
% inputs:       03-method/design.md · 03-method/figures/
% 验证方式:     与 design.md 逐条对照;框架图来自 scibox-diagram
% ============================================================
\section{Method}\label{sec:method}

\subsection{Problem Formulation}
<形式化:记号、输入输出、优化目标。>

\begin{equation}
  <目标函数/核心公式,与 design.md 编号一致>
  \label{eq:objective}
\end{equation}

\subsection{<核心模块 A>}
<机制描述;设计选择的理由,含"为什么不用 X"。>

\subsection{<核心模块 B>}
<同上。>

\subsection{Overall Framework}
% 框架图由 scibox-diagram 产出(draw.io 可编辑源文件在 03-method/figures/),
% 导出 PDF 放 figures/,此处引用:
% \begin{figure}[t]
%   \centering
%   \includegraphics[width=\linewidth]{framework.pdf}
%   \caption{<方法总览。读图即可理解主流程,不依赖正文。>}
%   \label{fig:framework}
% \end{figure}
<端到端流程走读。>

% 伪代码(可选):
% \begin{algorithm}
%   <design.md 中的伪代码,algorithm/algorithmic 宏包按需加>
% \end{algorithm}

% ============================================================
% 分节契约
% purpose:      用证据支撑 argument_map 的每个 claim
% allowed claims: 仅 evidence_table 已登记结果;auto 模式一律 \todobox 占位
% forbidden claims: 任何未经实验的数字以确定语气出现(铁律)
% inputs:       04-experiments/ · 05-ablation/ · experiment_log.md
% 验证方式:     nature-statistics 审查 + 数字可回溯 experiment_log
% ============================================================
\section{Experiments}\label{sec:exp}

\subsection{Setup}
\textbf{Datasets.}
<数据集、划分协议、评估指标;与 SOTA 对比须同协议,引 \cite{...}。>

\textbf{Baselines.}
<对比方法及来源; reproduced 数字标注出处。>

\textbf{Implementation details.}
<超参、种子控制、硬件;细节足以复现。>

\subsection{Main Results}
% 表格骨架就位;auto 模式数字全部占位,实验完成后回填并去掉 \todobox
\begin{table}[t]
  \centering
  \caption{<主结果:一句话说清该表支撑哪个 claim。均值$\pm$标准差,
  括号内种子数;统计规范过 nature-statistics。>}
  \label{tab:main}
  \begin{tabular}{lcc}
    \toprule
    Method & <Metric 1> & <Metric 2> \\
    \midrule
    <Baseline 1> \cite{...} & \todobox{待补实验} & \todobox{待补实验} \\
    <Baseline 2> \cite{...} & \todobox{待补实验} & \todobox{待补实验} \\
    \textbf{Ours}           & \todobox{待补实验} & \todobox{待补实验} \\
    \bottomrule
  \end{tabular}
\end{table}

<结果解读:指向表~\ref{tab:main},只做有证据的解读。>

\subsection{Ablation Study}\label{sec:ablation}
% 消融矩阵来自 05-ablation/ablation_plan.md;每个组件一行
\begin{table}[t]
  \centering
  \caption{<消融:去掉每个组件掉多少,为什么。>}
  \label{tab:ablation}
  \begin{tabular}{lc}
    \toprule
    Variant & <Metric 1> \\
    \midrule
    Full model              & \todobox{待补实验} \\
    w/o <组件 A>            & \todobox{待补实验} \\
    w/o <组件 B>            & \todobox{待补实验} \\
    \bottomrule
  \end{tabular}
\end{table}

<每行一句"去掉它掉多少,为什么"。审稿人可能问的
"为什么不消融 X"的预答也放这里。>

% ============================================================
% 分节契约
% purpose:      结果意味着什么、边界在哪、后续方向
% allowed claims: 从主结果直接推出的结论;局限须与范围边界一致
% forbidden claims: 超出证据的推广;防御性自我攻击(有礼有节地讲局限)
% inputs:       argument_map.md · 01-topic.md 范围边界
% 验证方式:     每个结论指向对应结果;anti-defensive-writing 过立场
% ============================================================
\section{Discussion}\label{sec:discussion}

<第 1 段:主结果回答了引言提出的问题到什么程度。>

<第 2 段:局限——与 topic.md 的范围边界一致,陈述事实不示弱。>

<第 3 段:后续方向。>

% Statistical analysis 小节(nature-statistics 审查后落笔):
% \subsection{Statistical analysis}
% <重复数、检验方法、多重比较校正、置信区间。>

% ============================================================
% 分节契约
% purpose:      收束:一句问题、一句方法、一句结果、一句意义
% forbidden claims: 新证据、新主张
% inputs:       各节成稿
% 验证方式:     与 abstract 逐句对齐,不引入新信息
% ============================================================
\section{Conclusion}\label{sec:conclusion}

<一句话问题回顾。><一句话方法。><一句话结果(\todobox{待补实验} 到位后回填)。>
<一句话意义与展望。>


\bibliographystyle{plainnat}
\bibliography{references}

\end{document}
